\section{Exact count generating function and the classical diffusion limit}
\label{app:finite-count}

The count component in \cref{prop:finite-count} is a critical linear
branching process with immigration after shifting by one. Its generating
function and diffusion are classical objects: the limiting diffusion is the
branching diffusion of \citet{Feller1951}, and the diffusion limit of
branching processes with immigration is due to \citet{KawazuWatanabe1971};
birth--death immigration approximations to Feller diffusions are discussed,
for example, by \citet{GiornoNobile2021}. We provide the specialization and
path argument to fix the factors in the present physical normalization.

\begin{proposition}[Generating function and count-scale limit]
\label{prop:count-diffusion}
Let $Y_t=L^n(t)-1$. For $y\in\{0,1,\ldots\}$ and $0\leq z\leq1$,
\begin{equation}
 \E[z^{Y_t}\mid Y_0=y]
 =\frac{u_t(z)^y}{1+bt(1-z)},\qquad
 u_t(z)=1-\frac{1-z}{1+bt(1-z)}.
 \label{eq:count-pgf}
\end{equation}
If $c_n=L_0^n/n\to c\geq0$, then, on every finite interval $[0,T]$,
\begin{equation}
 Z_n(\tau):=L^n(n\tau)/n\ \Longrightarrow\ Z(\tau)
 \quad\text{in }D([0,T],[0,\infty))\text{ with its }J_1\text{ topology},
 \label{eq:count-diffusion-limit}
\end{equation}
where $Z$ has continuous paths and solves, in the weak sense,
\begin{equation}
 \dd Z=b\dd\tau+\sqrt{2bZ}\dd W,\qquad Z(0)=c.
 \label{eq:count-sde}
\end{equation}
Its transition law is determined by
\begin{equation}
 \E_c e^{-\theta Z(\tau)}
 =\frac1{1+b\tau\theta}
   \exp\left(-\frac{c\theta}{1+b\tau\theta}\right),\qquad\theta\geq0,
 \label{eq:count-laplace}
\end{equation}
and
$\E_c Z(\tau)=c+b\tau$,
$\Var_c Z(\tau)=2bc\tau+b^2\tau^2$.
\end{proposition}

\begin{proof}
The shifted chain has birth rate $b(Y+1)$ and death rate $bY$. Realize it
as independent critical birth--death families, each with per-individual
birth and death rates $b$, together with immigration at Poisson rate $b$.
The generating function for one family started from one individual solves
$\partial_tu=b(u-1)^2$, $u_0(z)=z$, by conditioning on its first event.
Solving gives the displayed $u_t$. The $y$ initial families contribute
$u_t(z)^y$. Poisson immigration contributes
\[
 \exp\left(b\int_0^t[u_s(z)-1]\dd s\right)
 =[1+bt(1-z)]^{-1},
\]
proving \eqref{eq:count-pgf}. Nonexplosion from \cref{prop:finite-count}
justifies the family construction on finite intervals.

Substitute $t=n\tau$, $z=e^{-\theta/n}$, and $y=L_0^n-1$ into
\eqref{eq:count-pgf}; multiply by $e^{-\theta/n}$ to account for the shift.
Since $n(1-e^{-\theta/n})\to\theta$, the resulting transform converges
to \eqref{eq:count-laplace}. For fixed $\tau,\theta$ this convergence is
uniform when the scaled starting count stays in a bounded set. The limiting
transform is continuous in that starting count. The Markov property,
compact truncation of starting states, and this uniform transition
convergence identify all finite-dimensional laws. Equivalently one may
iterate the conditional Laplace transforms; products of exponentials form
a convergence-determining class for nonnegative finite-dimensional vectors.

To verify path tightness, the scaled martingale decomposition is
\begin{equation}
 Z_n(\tau)=c_n+b\tau+\mathcal M_n(\tau),\qquad
 \langle\mathcal M_n\rangle_\tau
 =\int_0^\tau(2bZ_n(s)-b/n)\dd s.
 \label{eq:scaled-count-martingale}
\end{equation}
Since $Z_n$ is a nonnegative submartingale,
\[
 \Prob\left(\sup_{\tau\leq T}Z_n(\tau)>R\right)
 \leq\frac{c_n+bT}{R}.
\]
This proves compact containment. Stop at
$\tau_R^n=\inf\{\tau:Z_n(\tau)>R\}$. The stopped bracket has rate at
most $2bR$ before this time. For stopping times $S_n\leq U_n\leq T$
with $U_n-S_n\leq\delta$, the martingale isometry gives
\[
 \E\bigl|\mathcal M_n(U_n\wedge\tau_R^n)
          -\mathcal M_n(S_n\wedge\tau_R^n)\bigr|^2
 \leq2bR\delta.
\]
The stopped drift changes by at most $b\delta$. Chebyshev's inequality,
followed by $\delta\downarrow0$ and then $R\uparrow\infty$, verifies the
stopping-time tightness criterion for $D$; see
\citet[Section 16]{Billingsley1999}. The jumps of $Z_n$ have size $1/n$,
so every subsequential limit is continuous. Finite-dimensional convergence
therefore identifies the whole path limit in \eqref{eq:count-diffusion-limit}.

For an explicit construction, let $B_1,B_2$ be independent standard
Brownian motions and set
\begin{equation}
 Z(\tau)=\frac b2
 \left[(\sqrt{2c/b}+B_1(\tau))^2+B_2(\tau)^2\right].
 \label{eq:count-bessel}
\end{equation}
It\^o's formula gives drift $b$ and martingale quadratic variation
$2bZ(\tau)\dd\tau$. Normalize the radial Brownian increment when $Z>0$
and choose a fixed unit direction when $Z=0$; its stochastic integral has
quadratic variation $\tau$ and hence is Brownian, giving
\eqref{eq:count-sde}. Rotational invariance of the two-dimensional Brownian
increments shows that \eqref{eq:count-bessel} is Markov and that its
transitions depend only on its current squared radius. A Gaussian integral
then gives \eqref{eq:count-laplace} with the current state in place of $c$.
These are precisely the limiting transitions already obtained. The same
Gaussian calculation gives the displayed mean and variance. Thus the limit
is a rescaled squared Bessel process of dimension two; it is independent
of daughter law and initial mass allocation once $b$ and $c$ are fixed.
\end{proof}

Order-one count fluctuations emerge on time $n$, whereas
\cref{thm:finite-discrepancy} already forces a mass-CDF discrepancy on a
logarithmic time scale. The classical diffusion is therefore a description
of a later count scale, not an explanation that requires count collapse at
the earlier mass-discrepancy horizon.
